\section{Waveform analysis framework}

This section summarizes the analysis path shared by the two final machine-learning models. The objective is to keep event definition, preprocessing, timing target, waveform windows, and blind-test evaluation fixed so that the main experimental variable is the model representation.

\subsection{Data partitioning}

The ROOT event population is split before any fitted event selection into a permanent development population and a permanent blind-test population. In the current default profile, \(40\%\) of the events are assigned to the blind-test split and \(60\%\) to development. The development population is then split with a validation fraction of \(20\%\), corresponding approximately to \(48\%\) training, \(12\%\) validation, and \(40\%\) blind test of the original population.

Fitted selection quantities are determined without using the blind-test population. The same persisted split identities are reused by the MLP and Onishi model studies, and the run-comparison utility checks the blind event identities before model-output correlations are computed.

\reportfigure
    {figures/data_split.pdf}
    {Dataset partitioning used by the final waveform pipeline. The permanent blind-test population is isolated before fitted selection and model development.}
    {fig:data-split}

\subsection{Event selection}

Event selection defines the coincidence population before timing analysis. The common energy-channel selection identifies the photopeak region for both detectors. For timing-channel studies, an additional time-over-threshold (ToT) selection isolates the main timing-pulse population. A baseline-noise requirement can be enabled by configuration when required.

All fitted selection cuts are derived from the development population and then frozen before application to the blind-test events.

\begin{table}[H]
    \centering
    \caption{Event-selection summary. Replace all placeholder entries with values exported from the reported run. Remove rows for disabled selection criteria.}
    \label{tab:event-selection}
    \begin{tabular}{lrrr}
        \toprule
        Selection stage & Input events & Remaining events & Efficiency [\%] \\
        \midrule
        Initial population & TODO & TODO & TODO \\
        Energy photopeak & TODO & TODO & TODO \\
        Pulse / trigger requirement & TODO & TODO & TODO \\
        Timing ToT & TODO & TODO & TODO \\
        Baseline RMS (optional) & TODO & TODO & TODO \\
        Final accepted population & TODO & TODO & TODO \\
        \bottomrule
    \end{tabular}
\end{table}


\reportfigure
    {figures/photopeak_selection.pdf}
    {Development-sample energy-photopeak selection.}
    {fig:photopeak-selection}

\reportfigure
    {figures/tot_selection.pdf}
    {Timing-channel ToT selection used in the timing-waveform analysis.}
    {fig:tot-selection}

\subsection{Native-time waveform preprocessing}

Only the waveform family required by the selected experiment mode is materialized. Waveforms are decoded and oriented consistently, clamped to detector-specific physical amplitude limits, and cropped around the main trigger while preserving the native acquisition-time grid. The current pipeline applies no denoising and no event-wise baseline subtraction.

For machine learning, amplitudes are mapped globally to \([0,1]\) using the configured detector-specific physical limits. This normalization is fixed by configuration rather than estimated from individual events or from the training population.

\reportfigure
    {figures/native_preprocessing_example.pdf}
    {Representative native-grid waveform and analysis region used for timing correction.}
    {fig:native-preprocessing}

\subsection{LED reference, calibration, and learning target}

Leading-edge discrimination (LED) provides the common timing reference. The waveform itself is not baseline-subtracted. Instead, for each event and detector, the configured LED threshold is interpreted as an offset above the mean pre-trigger baseline level measured in the configured baseline window. The threshold-crossing time is then obtained by linear interpolation on the rising edge. A fixed inter-channel timing offset is estimated from the training population:
\begin{equation}
    \widehat{C}_{12}
    =
    \operatorname{mean}_{\mathrm{train}}\!\left(\Delta t_{\mathrm{LED}}\right)
    -
    \mathrm{TOF}.
\end{equation}

The supervised target is the calibrated LED residual,
\begin{equation}
    y_{\mathrm{target}}
    =
    \Delta t_{\mathrm{LED}}
    -
    \mathrm{TOF}
    -
    \widehat{C}_{12}.
\end{equation}
The models therefore learn an event-by-event correction to the conventional timing estimate rather than predicting the physical TOF directly.

For final \texttt{model\_study} runs, the LED threshold offset above baseline is fixed in the experiment configuration and is shared by the MLP and Onishi studies being compared. LED-threshold dependence is investigated separately with the MLP-only \texttt{threshold\_scan} experiment; it is not folded into the architecture comparison.

\reportfigure
    {figures/led_threshold_selection.pdf}
    {Supporting LED-threshold study for the antisymmetric MLP. The final model comparison itself uses the fixed threshold declared by the paired model-study configurations.}
    {fig:led-threshold}

\subsection{Machine-learning input and waveform windows}

The two detector waveforms are synchronized to the selected LED reference and materialized on the native sampling grid. A temporal mask derived from the training split removes samples that are effectively constant in both detector channels. The mask is shared between the models for a given dataset and is then applied unchanged to validation and blind-test inputs.

The final profile defines two input windows:
\begin{align}
    \text{Onishi window:}\quad &[-1.5,\ 2.0]\ \mathrm{ns},\\
    \text{wide window:}\quad &[-2.0,\ 30.0]\ \mathrm{ns}.
\end{align}
The default temporal subsampling is one, so the retained samples remain on the native acquisition grid apart from the train-derived constant-sample mask.

\reportfigure
    {figures/ml_input_and_sample_mask.pdf}
    {Example detector-pair input and retained temporal samples after the training-derived constant-sample mask.}
    {fig:sample-mask}

\subsection{Final model-development protocol}

The two final models deliberately use different development policies because one is a proposed tunable model and the other is a fixed literature reference.

For the antisymmetric MLP, candidate hyperparameters are fitted using the training split only when more than one candidate is configured. An internal holdout carved from the training split is used for epoch-level early stopping, while the externally persisted validation split ranks the candidate configurations by validation CTR. After selection, the search checkpoint is discarded and a fresh final model is fitted on the complete development population. During this final fit, the MLP reserves the configured internal holdout from the development population for early stopping. If only one candidate is configured, validation-based model selection is skipped and the final development fit is performed directly.

The Onishi CNN exposes one fixed reference configuration and therefore performs no validation-based hyperparameter selection. Its final reference network is fitted once on the complete development population, i.e. training plus validation, and is then evaluated on the blind-test split. The temporal sample mask remains derived from the training split only.

This asymmetry in the fitting protocol must be stated explicitly when interpreting the model comparison: the Onishi model is treated as a fixed reference architecture, while the MLP uses validation only for hyperparameter selection when selection is actually required. Both final models are fitted on the development population before blind-test evaluation.

\subsection{CTR estimator}\label{sec:ctr-estimator}

Coincidence timing resolution (CTR) is evaluated from the timing-residual distribution~\cite{surti2015}. The repository uses the Gaussian-equivalent shortest empirical coverage interval. For a configured coverage fraction \(p\), the narrowest empirical interval containing \(\lceil pN\rceil\) finite residuals is rescaled to the equivalent Gaussian full width at half maximum. The current default coverage fraction is \(p=0.90\).

For a model prediction \(y_\theta\), the corrected residual is
\begin{equation}
    r_{\mathrm{ML}}=y_{\mathrm{target}}-y_\theta.
\end{equation}

Final CTR uncertainties are obtained from event bootstrap; the current default uses 500 bootstrap resamples. Presentation histograms are therefore only visual summaries and do not enter the CTR estimator.
